\chapter{What this is}
\label{ch:introduction}

Immersive Engineering is a technology mod built around the idea that machinery should look like
machinery. Its power is carried on wires strung between poles, its ore is crushed by a drum you can
watch turn, and the large machines are structures you assemble block by block rather than single
cubes that happen to contain a factory.

This manual documents \textbf{LDImmersiveEngineering}, a private fork of that mod for Minecraft
1.12.2. Everything the original does, this does. On top of it sit four large additions and one
performance regime, none of which exist upstream:

\begin{description}
\item[The virtual power grid] Moves Immersive Flux between places with no wire between them.
      \autoref{ch:grid}.
\item[Petroleum] A complete oil chain: prospect, drill, refine, store, burn, sell.
      \autoref{ch:petroleum}.
\item[The virtual fluid network] The same idea as the grid, for fluid. \autoref{ch:fluidnet}.
\item[Conduits] Surface-mounted bundled wiring for the inside of buildings. \autoref{ch:conduits}.
\item[City mode] Trades simulation detail for server tick time, subsystem by subsystem.
      \autoref{ch:citymode}.
\end{description}

\section{How to read this}

Part~II is Immersive Engineering as it ships. Part~III is what the fork adds. Part~IV is generated
straight from the source tree --- every block, every recipe, every configuration option --- and is
the place to look when you want a number rather than an explanation.

Three kinds of callout appear throughout:

\begin{ietip}
Something that makes the mod easier, or a habit worth forming early.
\end{ietip}

\begin{iewarning}
Something that will cost you a build, a machine or an afternoon if you get it wrong.
\end{iewarning}

\begin{iefork}
Behaviour that differs from stock Immersive Engineering. If you already know the mod, these boxes
are the parts you have to read.
\end{iefork}

\section{Where the numbers come from}

Every figure in this manual was read out of the source. Where a number is configurable --- and most
are --- the value quoted is the shipped default, and \autoref{ch:config} lists the option that
changes it.

\begin{iewarning}
A modpack may have changed any of them. If a machine in your world does not match a number here,
check the config before assuming the manual is wrong.
\end{iewarning}
